% 五种结构不同的课程报告版式。内容与字段由 common/metadata/body 提供。
% 使用版本化文件，避免覆盖各实验已安装或手工修改的 common.tex。
\raggedbottom
\setlength{\columnseprule}{0pt}
\definecolor{ReportTint}{HTML}{F1F4F7}
\newcommand{\ReportInfoRows}{%
  作者（姓名）：\ReportAuthor\quad 学号：\ReportStudentID\\[4pt]
  指导老师：\ReportTeacher\quad 实验时间：\ReportExperimentDate\par}
\newcommand{\ReportDraftNotice}{%
  \ifReportDraft
    \noindent\fcolorbox{black!25}{black!3}{\parbox{\dimexpr\linewidth-2\fboxsep-2\fboxrule\relax}{%
      \small 模板草稿：所有“待补充”均须以讲义、原始记录或可核实资料填写；占位图表与文献不能作为实测结果提交。
    }}\par\vspace{6pt}
  \fi}
\newcommand{\ReportChineseSummary}{%
  \noindent\textbf{摘要}\quad\ReportAbstractCN\par
  \noindent\textbf{关键词：}\ReportKeywordsCN\par}
\newcommand{\ReportEnglishSummary}{%
  \noindent\textbf{Abstract}\quad\ReportAbstractEN\par
  \noindent\textbf{Keywords:}\ \ReportKeywordsEN\par}
\newcommand{\ReportEnglishTitle}{%
  {\bfseries\ReportTitleEN\par}\vspace{3pt}
  {\small\ReportAuthorEN\par}}
\newcommand{\ReportClassicFront}{%
  \vspace*{8mm}
  {\centering\ReportTitleFont\ReportTitle\par}
  \begin{center}\ReportInfoRows\end{center}
  \ReportDraftNotice
  \ReportChineseSummary\vspace{7pt}
  {\centering\ReportEnglishTitle}\vspace{5pt}
  \ReportEnglishSummary
  \par\vspace{7pt}\hrule height 0.6pt\vspace{2pt}\hrule\vspace{8pt}}
\newcommand{\ReportJournalFront}{%
  \vspace*{8mm}
  {\centering\ReportTitleFont\ReportTitle\par\vspace{5pt}\ReportEnglishTitle}
  \begin{center}\small\ReportInfoRows\end{center}
  \ReportDraftNotice
  \noindent\begin{minipage}[t]{0.475\textwidth}\small\ReportChineseSummary\end{minipage}%
  \hfill\begin{minipage}[t]{0.475\textwidth}\small\ReportEnglishSummary\end{minipage}
  \par\vspace{8pt}\hrule\vspace{8pt}}
\newcommand{\ReportEditorialFront}{%
  \vspace*{8mm}
  {\raggedright\ReportTitleFont\ReportTitle\par}\vspace{7pt}
  {\raggedright\ReportEnglishTitle}\vspace{8pt}
  {\small\noindent\ReportInfoRows}\vspace{8pt}
  \ReportDraftNotice
  \begingroup\setlength{\fboxsep}{9pt}
  \noindent\colorbox{ReportTint}{\parbox{\dimexpr\textwidth-2\fboxsep\relax}{%
    \small\ReportChineseSummary}}\par\endgroup\vspace{7pt}
  {\small\ReportEnglishSummary}\vspace{10pt}
  {\color{ReportAccent}\hrule height 1.2pt}\vspace{8pt}}
\newcommand{\ReportBriefFront}{%
  \vspace*{8mm}
  \begingroup\setlength{\fboxsep}{10pt}
  \noindent\colorbox{ReportTint}{\parbox{\dimexpr\textwidth-2\fboxsep\relax}{%
    {\small\heiti\color{ReportAccent}物理实验报告\par}\vspace{5pt}
    {\raggedright\ReportTitleFont\ReportTitle\par}\vspace{5pt}
    \ReportEnglishTitle}}\par\endgroup
  \begin{center}\small\ReportInfoRows\end{center}
  \ReportDraftNotice
  \ReportChineseSummary\vspace{5pt}
  {\small\ReportEnglishSummary}\vspace{7pt}
  {\color{ReportAccent}\hrule height 1pt}\vspace{7pt}}
\newcommand{\ReportLongFront}{%
  \vspace*{2cm}
  {\centering\small\heiti 物理实验报告\par}\vspace{1cm}
  {\centering\ReportTitleFont\ReportTitle\par}\vspace{10pt}
  {\centering\ReportEnglishTitle}\vspace{1cm}
  \begin{center}\ReportInfoRows\end{center}\vspace{1cm}
  \ReportDraftNotice
  {\small\ReportChineseSummary}\vspace{10pt}
  {\small\ReportEnglishSummary}
  \thispagestyle{plain}\clearpage}
\ifcase\ReportTemplateID\relax
\or % 01 经典课程：双栏、居中题头、上下双语摘要。
  \renewcommand{\ReportTitleFont}{\heiti\fontsize{19pt}{25pt}\selectfont}
  \renewcommand{\baselinestretch}{1.15}
  \setlength{\columnsep}{8mm}
\or % 02 紧凑学术：双栏、并排摘要、细栏间线。
  \renewcommand{\ReportTitleFont}{\songti\bfseries\fontsize{17pt}{22pt}\selectfont}
  \renewcommand{\baselinestretch}{1.06}
  \setlength{\columnsep}{8mm}
  \setlength{\columnseprule}{0.25pt}
  \ctexset{section={format=\normalsize\heiti,beforeskip=10pt,afterskip=4pt}}
\or % 03 科研杂志：单栏引言，其后同页进入双栏。
  \renewcommand{\ReportTitleFont}{\heiti\fontsize{25pt}{31pt}\selectfont}
  \renewcommand{\baselinestretch}{1.12}
  \definecolor{ReportAccent}{HTML}{174C68}
  \definecolor{ReportTint}{HTML}{EEF5F7}
  \setlength{\columnsep}{9mm}
  \ctexset{section={format=\large\heiti\color{ReportAccent}}}
\or % 04 技术简报：单栏、灰底题头、无首行缩进、段落留白。
  \renewcommand{\ReportTitleFont}{\heiti\fontsize{19pt}{25pt}\selectfont}
  \renewcommand{\baselinestretch}{1.08}
  \definecolor{ReportAccent}{HTML}{225A50}
  \definecolor{ReportTint}{HTML}{EFF4F1}
  \setlength{\parindent}{0pt}
  \setlength{\parskip}{4pt}
  \ctexset{section={format=\large\heiti\color{ReportAccent},beforeskip=10pt,afterskip=4pt}}
\or % 05 学术长报告：独立题头页、单栏正文、较大字号与行距。
  \renewcommand{\ReportTitleFont}{\songti\bfseries\fontsize{25pt}{32pt}\selectfont}
  \renewcommand{\baselinestretch}{1.18}
  \ctexset{section={format=\Large\heiti,beforeskip=12pt,afterskip=6pt}}
\fi
\makeatletter
\renewcommand{\MakeReportFrontMatter}{%
  \ifcase\ReportTemplateID\relax
  \or\twocolumn[{\begin{@twocolumnfalse}\ReportClassicFront\end{@twocolumnfalse}}]
  \or\twocolumn[{\begin{@twocolumnfalse}\ReportJournalFront\end{@twocolumnfalse}}]
  \or\ReportEditorialFront
  \or\ReportBriefFront
  \or\ReportLongFront
  \fi
  \ifnum\ReportTemplateID<5\thispagestyle{plain}\fi}
\makeatother
\newif\ifReportColumnsStarted
\newif\ifReportColumnsOpen
\newcommand{\ReportEndColumns}{%
  \ifReportColumnsOpen\end{multicols}\ReportColumnsOpenfalse\fi}
\ifnum\ReportTemplateID=3
  \usepackage{multicol}
  \AddToHook{cmd/section/before}{%
    \ifReportColumnsStarted\else
      \ifnum\value{section}=1\relax
        \ReportColumnsStartedtrue\ReportColumnsOpentrue
        \begin{multicols}{2}\raggedcolumns
      \fi
    \fi}
  \AddToHook{cmd/appendix/before}{\ReportEndColumns}
  \AtEndDocument{\ReportEndColumns}
\fi
